\documentclass[conference,a4paper]{IEEEtran}
\usepackage[T1]{fontenc}
\usepackage{newtxtext,newtxmath}
\usepackage{graphicx,booktabs,amsmath,multirow,array}
\usepackage{microtype}
\usepackage[hidelinks]{hyperref}
\usepackage{xurl}
\renewcommand{\arraystretch}{1.12}
\newcommand{\amber}{\textsc{amber}}
\newcommand{\red}{\textsc{red}}
\newcommand{\watch}{\textsc{watch}}

\begin{document}

\title{Countersign: On-Device Detection of Impersonation Scam Calls with Family-Linked Response}

\author{\IEEEauthorblockN{Pushkar Joshi, Sarah Mittal, Reshaan Nadar}
\IEEEauthorblockA{Team \emph{losers}, RAKSHAM 2026 Finale, Track 1 (voice-cloning fraud)\\
Code, data cards and every result: \url{https://github.com/dope-XD/countersign}}}

\maketitle

\begin{abstract}
Impersonation calls (``digital arrest'', fake bank KYC, a relative in trouble) persuade victims to transfer money
while the call is still in progress. We present Countersign, an Android application that detects such calls on the
victim's own phone and brings in the victim's family before money moves. Speech is transcribed on the device, and
two independent judges assess the conversation: a multilingual encoder (e5-small) fine-tuned on 267{,}601
conversation windows, which scores the call as a whole, and a negation-aware rule base of 19 scam ``moves'' that
supplies the explanation shown to the user. A warning is raised when either judge is confident; the family alarm
requires two independent sources. A separate check asks a family member's own phone whether it placed the call,
which defeats voice cloning and caller-ID spoofing without using audio. On 1{,}019 genuine bank conversations never
used in training, 2.3\% raise a warning (a keyword baseline: 50.1\%); 87.8\% of 4{,}499 held-out scam conversations
do. We report the weaknesses we measured, among them scripted English scams and talk about scams.
\end{abstract}

\begin{IEEEkeywords}
scam detection, social engineering, on-device inference, speech, multilingual, privacy
\end{IEEEkeywords}

\section{Introduction}

In a typical ``digital arrest'' call, a caller claiming to be from the CBI tells the victim that a parcel in her
name contained drugs, keeps her on the line for an hour, forbids her to tell anyone, and finally asks her to move her
savings to a ``safe account'' to clear her name. Fake KYC updates and ``your son has had an accident'' calls follow
the same script. Two properties make these calls hard to stop: the harm happens \emph{during} the call, and the
victim is deliberately isolated from the people who would recognise the scam.

A detector for this setting must therefore work in real time, on the victim's phone, in the languages the call is
held in (English, Hindi and their mixture, Hinglish), and it must hand the decision to a person the victim trusts.
It must also not alarm on genuine calls: a real bank's fraud team uses the same vocabulary (\emph{OTP},
\emph{account}, \emph{police}) as the scammer. Our first prototype, which scored keyword cues, warned on half of
genuine bank calls (Table~\ref{tab:icfd}); this observation shaped the design below.

\textbf{Contributions.}
\begin{itemize}
\item An on-device detection pipeline that judges \emph{what the caller is doing} across the whole conversation,
using a fine-tuned multilingual encoder, combined with interpretable cue rules through a two-source rule for the
highest alert level (Section~\ref{sec:detect}).
\item A caller check that needs no audio: the claimed family member's phone reports its own calls, so a cloned voice
or a spoofed number fails because that phone never placed the call (Section~\ref{sec:response}).
\item A response path that reaches the family (full-screen alarm, SMS fallback, escalation) and a two-tap recovery
from false alarms, with privacy constraints enforced in the architecture (Sections~\ref{sec:response}, \ref{sec:privacy}).
\item An evaluation on data the models never saw, with false-positive rates and named failure cases
(Section~\ref{sec:eval}); all numbers are reproducible from the repository.
\end{itemize}

\section{Related Work}

\textbf{Deployed on-device detectors.} Google's Scam Detection analyses calls on Pixel phones for ``patterns commonly
associated with scams'' and keeps no audio or transcript~\cite{google}. It is English-only, limited to specific
handsets, and does not involve the victim's family.

\textbf{Models for conversational scams.} PreScam~\cite{prescam} collects 11{,}573 scam conversations from user
reports and finds that a fine-tuned BERT predicts scam progression better than zero-shot large language models
(83.4 vs.\ 63.3--69.9 AUC). SCRIPTMIND~\cite{scriptmind} models scams as crime scripts with fine-tuned small
language models, and Shen et~al.~\cite{shen} study real-time LLM warnings and the trade-off between timeliness and
precision. These results motivate our choice of a compact fine-tuned encoder over an on-phone generative model.

\textbf{Data.} ICFD-31k~\cite{icfd} provides about 31{,}000 bilingual (English and Hinglish) bank-call conversations,
synthetic and generated with Llama~3.3~70B, with streaming verdicts every 3~s, ten fraud families in the main set and
five unseen fraud types in a separate cross-domain set. It is our main training and test source.

\textbf{Synthetic voice detection.} Leaderboards of speech deepfake detectors report that small models trade accuracy
for efficiency~\cite{dfarena}. We use the first layers of a post-trained AntiDeepfake model~\cite{antideepfake} and,
given its error rate over phone audio (Section~\ref{sec:eval}), treat it as weak evidence only.

\section{System Overview}

\begin{figure*}[t]
\centering
\includegraphics[width=\textwidth]{fig/architecture.pdf}
\caption{Architecture. All listening and judging happen on the protected user's phone. The relay carries identifiers
only; the detection server is optional and receives audio only after the phone itself has raised a warning.}
\label{fig:arch}
\end{figure*}

Fig.~\ref{fig:arch} shows the three parties. Every family member installs the app. The \emph{protected user's
phone} captures the call, transcribes it and decides; the \emph{relay}, a small server without models, links family
phones and carries call events, caller checks and alerts as identifiers; an optional \emph{detection server} gives a
second opinion on calls the phone already finds suspicious.

\textbf{Threat model.} The attacker calls from an unknown number or spoofs a known one, impersonates an authority, a
bank or a family member, may use a cloned voice, and steers the victim towards a transfer, a one-time password, or
installing a remote-access app. We assume the attacker cannot control the family member's phone. Detection requires
the call to be on speaker (Section~\ref{sec:limits}).

\textbf{Alert levels.} A call is at one of four levels: \emph{none}, \watch{} (a few signs; no visible warning),
\amber{} (a warning card naming the scam and the reasons) and \red{} (``This is likely a scam. Hang up.'' and the
family alarm). Within a call the level never decreases, except when the user declares the call genuine.

\section{On-Device Detection}
\label{sec:detect}

\subsection{Audio capture and speech recognition}
Android lets an accessibility service keep the microphone during a call; the app uses one to record 16~kHz mono
audio into memory while a call is active and listening is enabled. Silero VAD~\cite{silero} segments speech
(threshold 0.5, minimum silence 0.4~s, segments up to 12~s). Each segment is transcribed by Meta's Omnilingual
ASR~\cite{omniasr} (300M parameters, CTC, int8) through sherpa-onnx, and its audio is discarded. We selected this
recogniser on our 59 spoken calls: its word error rate was 28\%, against 65\% and 44\% for the Hindi and English
IndicConformer models, each of which also loses the other language of a code-mixed call.

\subsection{Judge 1: the conversation model}
\begin{figure*}[t]
\centering
\includegraphics[width=\textwidth]{fig/decision.pdf}
\caption{Decision. Both judges see the transcript and the phone's context. \amber{} follows the more confident
judge; \red{} requires two independent sources.}
\label{fig:decision}
\end{figure*}

The first judge reads the conversation so far and estimates whether the caller has made a scam move.

\emph{Model.} We fine-tune \texttt{multilingual-e5-small}~\cite{e5} (118M parameters) with mean pooling and a logistic
head. Inputs are the most recent part of the call (left-truncated to 320 tokens, which covers 99\% of 1{,}200-character
windows), normalised but kept in its original script, with the e5 ``query:'' prefix.

\emph{Training data.} ICFD-31k streaming windows plus our own sets (Hinglish call lines, public scam dialogues and SMS
corpora, and template-generated genuine calls that use scam vocabulary), 267{,}601 windows in total with 24.6\%
positive. Our own rows are repeated eight times to balance the synthetic majority. We train for one epoch with AdamW and a linear schedule with warm-up, and select the checkpoint on the
development set only (mean per-source AUC 0.989).

\emph{Splits.} ICFD's validation and test splits contain only fraud, so we hold out conversations of its training
split by a hash of the conversation id: 10\% for development and 10\% as the only ICFD test set with genuine calls.
ICFD's own test split and its cross-domain split stay untouched.

\emph{On the phone.} The model runs as an fp32 ONNX graph (470~MB) with onnxruntime; it reproduces PyTorch to within
$10^{-5}$, whereas int8 quantisation shifted scores by up to 0.20 and was rejected. Tokenisation uses our own Kotlin
implementation of the SentencePiece unigram model, identical to the reference tokenizer on 3{,}005 test texts.

\emph{Accumulation.} Let $p_t$ be the model's score after segment $t$ and $s_t=\tfrac12(p_t+p_{t-1})$. The call reaches
\begin{equation}
\mathrm{level}_t=\begin{cases}
\red{} & \text{if } s_t \ge 0.9855 \text{ and } s_{t-1} \ge 0.9855,\\
\amber{} & \text{if } s_t \ge 0.142,\\
\watch{} & \text{if } s_t \ge 0.0102,\\
\text{none} & \text{otherwise,}
\end{cases}
\end{equation}
and the level of the call is the highest $\mathrm{level}_t$ reached so far.
The thresholds were set on 1{,}788 genuine development calls at 10\%, 3\% and 0.5\% false-positive rate. A known
contact or a verified caller raises them (0.142, 0.9865, 0.9932). A failed caller check or an opened remote-access
app is sufficient for \amber{} on its own. Together with \amber{} from speech, these signals, an unknown video
call or an international number produce \red{}.

\emph{Reasons.} A move matcher embeds each segment with the base e5 model and compares it with written descriptions
of 17 scam moves (e.g.\ ``claims to be police'', ``asks to keep it secret'', ``asks for the OTP''). Each move has its
own threshold, the 99.5th percentile of similarity on held-out genuine turns. A move is reported only when $p_t$ rose
by at least 0.15 or the level changed. The matcher never sets the level.

\subsection{Judge 2: cue rules}
The second judge is a rule base of phrases in English, Hinglish and Hindi grouped into 19 categories (authority,
threat, secrecy, money request, credentials, remote access, urgency, and others). A negation window prevents
``we will never ask for your OTP'' from counting as a request; a \emph{discussion} category (talk about scams) never
counts towards \red{}. Combinations of categories name the scam type (e.g.\ \emph{digital arrest} = authority plus
a threat, secrecy or a video-call demand). The rules reach \watch{} with one category, \amber{} with two strong
categories, and \red{} with four including an action and a type. A context signal counts once: it is never also used
as a cue towards the same \amber{}.

\subsection{Context and voice evidence}
The phone contributes signals of its own: an unknown or international number, a remote-access app (AnyDesk,
TeamViewer and others) or a payment app brought to the front during the call, a call longer than 15 minutes, a failed
caller check, the user reading out four to eight digits on an unknown call, and an unknown number claiming to be a
bank or authority. A voice-clone check scores the first 20~s of speech with the first six layers of AntiDeepfake
wav2vec-large~\cite{antideepfake} and a logistic head (int8, 127~MB). Because of its error rate over phone audio
(Table~\ref{tab:voice}), a synthetic-voice signal can raise \watch{} to \amber{} but never produces \red{}.

\subsection{Combination}
Fig.~\ref{fig:decision} summarises the rule. \amber{} is the higher of the two judges, except that cue words alone
do not produce \amber{} while the conversation model is below \watch{}: genuine bank and courier calls use the same
words. \red{} requires two independent sources: one judge at \red{} from speech together with the other judge at
\amber{} or above; a judge that already paired speech with a strong phone signal; the user pressing ``Alert family'';
or the detection server. The server, when present, can raise the level but never lower it.

\subsection{Optional second opinion}
Once the phone itself is at \amber{}, it may stream the call to the detection server, which runs a larger recogniser
(Omnilingual 1B), the same conversation model, and the full 24-layer AntiDeepfake model on 4~s windows. Audio is
kept in a rolling 20~s buffer in memory and dropped when the stream closes.

\section{Caller Check and Response}
\label{sec:response}

\begin{figure*}[t]
\centering
\includegraphics[width=0.94\textwidth]{fig/sequence.pdf}
\caption{Message flow for one scam call: caller check, optional second opinion, alert, family alarm with SMS
fallback and escalation, and the false-alarm path. Dashed: optional or fallback paths.}
\label{fig:seq}
\end{figure*}

\textbf{Caller check.} When a caller claims to be a family member, or calls from a family member's number, the relay
compares the claim with that member's own reported call events around the moment the call rang
(Fig.~\ref{fig:seq}). An outgoing call to the protected user verifies the caller; no call while the member's app was
seen in the last two minutes yields ``Not Rahul's phone''; an app not seen for two minutes yields ``could not
check'', never a false negative. Negative verdicts wait up to 6~s for late events. Numbers are matched to member
identifiers on the phone and never leave it. The user may also ring the member, whose phone asks ``Is it you?''; a
``yes'' is signed by a hardware-backed key after biometric confirmation.

\textbf{Alert and escalation.} The first time a call reaches \red{}, the phone sends one alert carrying the call
identifier, level, scam type and cue identifiers. The relay rings every family member whose app is online with a
full-screen alarm (``Mummy may be on a scam call'', with \emph{Call} and \emph{Message} actions) and returns the
offline members, whom the phone texts from its own SIM. If nobody acknowledges within 60~s, the phone texts all
members and places a call to the first one.

\textbf{Recovering from false alarms.} ``I know this caller'' or ``It's not a scam'', confirmed once, stops warnings
for the call, closes any server stream, clears the alarm on family phones (``False alarm: Mummy says it's fine'')
and stores one identifier-only record for human review.

\section{Privacy}
\label{sec:privacy}

\begin{table}[t]
\centering
\caption{Where each kind of data goes}
\label{tab:privacy}
\small
\begin{tabular}{@{}p{0.27\columnwidth}p{0.66\columnwidth}@{}}
\toprule
Data & Handling \\
\midrule
Call audio & Phone memory, discarded per segment; to the server only after \amber{}, in memory; one 20~s AES-256 encrypted clip for a confirmed scam, deleted after 30 days \\
Transcript & Phone memory for the duration of the call; never stored or sent \\
Phone numbers & On the phone only; matched to member identifiers by the last ten digits \\
Relay & Call events (member ids), check verdicts, alert (level, type, cue ids) \\
Logs & Levels, types and cue identifiers; never words or numbers \\
\bottomrule
\end{tabular}
\end{table}

Table~\ref{tab:privacy} lists every data flow. Detection, the explanation and the caller check work without sending
speech or text anywhere; the relay cannot reconstruct what was said. ``Check a recording or message'' applies the same
analyser to a file or pasted text, entirely on the phone.

\section{Implementation}

The Android client is written in Kotlin; the server, the training code and the evaluation are in Python. Every
decision component exists in both languages, and continuous integration checks that both give identical results on
shared test vectors (tokenizer, conversation accumulator, cue rules, combination). The model pack is a single
798~MB download, verified by hash, sized for a phone with 6~GB of RAM; detection works offline after the download. Of the
product code, 66\% was written during the 48-hour finale (measured with \texttt{git blame} against the check-in
commit).

\section{Evaluation}
\label{sec:eval}

\textbf{Protocol.} Thresholds were chosen on development data only. Test sets were never used for training,
threshold selection or checkpoint choice. A call counts as \emph{flagged} when it reaches \amber{} or \red{} at any
point. ICFD conversations are fed in their streaming chunks, our calls segment by segment as the phone does.

\begin{table}[t]
\centering
\caption{ICFD-31k conversations not used in training: share reaching \amber{} / \red{}}
\label{tab:icfd}
\small
\begin{tabular}{@{}lrcc@{}}
\toprule
Set & $n$ & Countersign & Keyword baseline \\
\midrule
Genuine, clear & 1{,}019 & 2.3 / 0.4 & 50.1 / 8.9 \\
Genuine, ambiguous & 511 & 4.5 / 0.8 & 70.8 / 15.5 \\
\midrule
Scam, test split & 4{,}499 & 87.8 / 59.9 & 94.3 / 51.5 \\
Scam, held out (clear) & 193 & 99.5 / 82.9 & -- \\
Scam, held out (subtle) & 370 & 98.9 / 75.9 & -- \\
Unseen types (clear) & 760 & 75.8 / 42.9 & 90.9 / 34.7 \\
Unseen types (subtle) & 240 & 54.6 / 17.5 & 78.3 / 22.5 \\
\bottomrule
\end{tabular}

\vspace{2pt}\footnotesize Percent. Median time to \amber{}: 45~s (test split), 36~s (held out).
\end{table}

\textbf{False alarms on genuine calls.} Table~\ref{tab:icfd} shows the effect of judging the conversation instead of
keywords: on clear genuine bank calls, warnings fall from 50.1\% to 2.3\% and alarms from 8.9\% to 0.4\%. The keyword
baseline catches more scams only because it flags most genuine calls as well.

\textbf{Scams.} 87.8\% of scam conversations in ICFD's test split reach \amber{} and 59.9\% reach \red{}, half a
minute to a minute into the call. Detection drops on the five fraud types absent from training (75.8\% clear,
54.6\% subtle), which bounds what the model generalises.

\begin{table}[t]
\centering
\caption{Other held-out sets (text, fed turn by turn)}
\label{tab:ours}
\small
\begin{tabular}{@{}lcccc@{}}
\toprule
Set & \multicolumn{2}{c}{$n$ (scam / genuine)} & Scams & Genuine \\
\midrule
Team calls & 16 & 13 & 93.8\% & 7.7\% \\
Adversarial rewrites & 16 & 14 & 75.0\% & 21.4\% \\
Hinglish call lines & 124 & 22 & 100\% & 0\% \\
Near-domain genuine & 0 & 30 & -- & 10.0\% \\
Fresh set$^\dagger$ & 10 & 10 & 10/10 & 2/10 \\
\midrule
Spot-check kit & \multicolumn{4}{l}{28 of 30 items correct} \\
\bottomrule
\end{tabular}

\vspace{2pt}\footnotesize Share flagged (\amber{} or \red{}). $^\dagger$Calls over a simulated phone line.
\end{table}

\textbf{Our sets and unseen inputs.} Table~\ref{tab:ours} reports the combined system on our own held-out sets. The
fresh set was written after the system was frozen and used for nothing else: six calls spoken by voices absent from
all training data, each also passed through a simulated phone line, and four messages per class. All ten scams were
flagged. The genuine items flagged were people discussing a scam and a bank fraud alert read by a synthetic voice.

\textbf{Audio end to end.} On 59 spoken calls (32 scams) passed through the recogniser, the conversation model alone
flags 66\% of scams on clean audio, 62\% over G.711 and 50\% after loudspeaker replay plus G.711, with 7--11\% of
genuine calls at \amber{} and at most 4\% at \red{}. Here the keyword path did better on scams (88\% over G.711):
our English scam scripts are more formal than ICFD's bank-call English. The combined system has not yet been measured
on this audio set.

\begin{table}[t]
\centering
\caption{Voice-clone check (shipped int8 model)}
\label{tab:voice}
\small
\begin{tabular}{@{}lcc@{}}
\toprule
Held-out set & Channel & Equal error rate \\
\midrule
In-the-Wild & clean / G.711 & 23.2 / 24.6\% \\
ReplayDF (replayed by speaker) & G.711 & 35.2\% \\
Stock TTS vs.\ real Hindi & clean / G.711 & 10.5 / 17.5\% \\
\bottomrule
\end{tabular}
\end{table}

\textbf{Voice.} Table~\ref{tab:voice} explains why the voice check is weak evidence: equal error rates of 23--35\%
on realistic channels are too high to justify an alarm on their own.

\section{Limitations and Future Work}
\label{sec:limits}

\textbf{Measured weaknesses.} (i) Scams read politely from a script in English can stay below \amber{}. (ii) People
talking \emph{about} a scam can trigger \amber{}, occasionally \red{}. (iii) A genuine call in a synthetic voice can
reach \amber{}; the card then names the voice as the reason. (iv) Short calls rarely reach \red{}, since it requires
evidence over two segments. (v) Recognition errors over a loudspeaker lower detection (Section~\ref{sec:eval}).

\textbf{Deployment constraints.} On the earpiece the microphone mostly hears the user, so the call must be on
speaker, and some handsets' echo cancellation removes the caller even then. ICFD-31k is synthetic; performance on
real recorded calls needs a pilot with consented data.

\textbf{Not yet built.} A full-screen block over payment apps (designed; today the call card and home screen warn), a
small on-device language model as a third judge (prompt and parser written and tested, no runtime yet), three
specified phone signals (screen sharing, long unknown video calls, 1600-series numbers), and the combined system on
the server.

\section{Conclusion}

Countersign shows that impersonation scams can be detected on the victim's phone, in Indian languages, with an
explanation the victim can read, and that the decisive step, bringing in the family, can be automated without
sending the conversation anywhere. Judging the whole conversation rather than its keywords cut false warnings on
genuine bank calls from one in two to about two in a hundred.

\begin{thebibliography}{10}
\bibitem{prescam} W.~Sun, S.~Ma, Y.~Li, T.~Ma, Z.~Wang, C.~Nelson, X.~Xiao, and Y.~Ye, ``PreScam: A benchmark for predicting scam progression from early conversations,'' arXiv:2605.12243, 2026.
\bibitem{icfd} R.~Ahuja, K.~Prateek, and S.~Singh, ``ICFD-31k: A large-scale dataset and benchmark for real-time conversational fraud detection,'' in \emph{Proc.\ IJCAI}, 2026, pp.~7065--7073.
\bibitem{scriptmind} H.~Kim, C.~Kim, S.~Shin, and J.~Kang, ``SCRIPTMIND: Crime script inference and cognitive evaluation for LLM-based social engineering scam detection system,'' arXiv:2601.13581, 2026.
\bibitem{shen} Z.~Shen, S.~Yan, Y.~Zhang, X.~Luo, G.~Ngai, and E.~Y.~Fu, ``\,`It warned me just at the right moment': Exploring LLM-based real-time detection of phone scams,'' arXiv:2502.03964, 2025.
\bibitem{e5} L.~Wang, N.~Yang, X.~Huang, L.~Yang, R.~Majumder, and F.~Wei, ``Multilingual E5 text embeddings: A technical report,'' arXiv:2402.05672, 2024.
\bibitem{omniasr} Omnilingual ASR Team (Meta FAIR), ``Omnilingual ASR: Open-source multilingual speech recognition for 1600+ languages,'' arXiv:2511.09690, 2025.
\bibitem{antideepfake} W.~Ge, X.~Wang, X.~Liu, and J.~Yamagishi, ``Post-training for deepfake speech detection,'' arXiv:2506.21090, 2025.
\bibitem{dfarena} S.~Dowerah \emph{et al.}, ``Speech DF Arena: A leaderboard for speech deepfake detection models,'' arXiv:2509.02859, 2025.
\bibitem{silero} Silero Team, ``Silero VAD: Pre-trained enterprise-grade voice activity detector,'' GitHub, 2024.
\bibitem{google} Google, ``New real-time protections on Android,'' Google Security Blog, Nov.~2024.
\end{thebibliography}

\end{document}
